\documentclass[11pt]{article}
\usepackage[margin=1in]{geometry}
\usepackage{amsmath,amssymb,amsthm,booktabs}
\usepackage[hidelinks]{hyperref}
\setlength{\emergencystretch}{2em}
\newtheorem{theorem}{Theorem}
\newtheorem{proposition}[theorem]{Proposition}
\newtheorem{corollary}[theorem]{Corollary}
\newcommand{\C}{\mathbb C}
\newcommand{\Z}{\mathbb Z}
\newcommand{\im}{\operatorname{im}}
\newcommand{\GL}{\operatorname{GL}}
\title{Meridian-varying tangents and formal stabilization at icosahedral points}
\author{Internal research working proof}
\date{16 November 2026}
\begin{document}
\maketitle

All sixteen retained $A_5$ quotient-kernel classes, supported on
twelve of the sixteen target knots, give complex three-dimensional
$H_3$ reflection representations. Their meridian-varying KCH
tangent spaces have now been computed exactly. After normalizing
one meridian projector, each has dimension five: four residual
conjugations and one direction with nonzero meridian-eigenvalue
derivative. The fixed-eigenvalue formal representations are therefore
conjugate to the original representation, strengthening the earlier
formal-orthogonality boundary.

A separate uniform splitting argument gives an all-order restriction
even when the eigenvalue varies: adjoining any finite-dimensional
trivial summand cannot deform these points into irreducible
higher-dimensional KCH representations. Thus this particular route
cannot raise the rank-three lower bound to four. It does not exclude
other representations or components.

The surviving rank-three tangent is not an actual curve certificate.
Both embeddings modulo $11$ pass through degree three, with a separate
polynomial replay, but characteristic-zero integrability remains open.
No target knot is settled. The uniform arguments are written working
proofs, unformalized and externally unreviewed, without a priority claim.

\section{Exact normalized tangent spaces}

Let $\Gamma$ be a knot group and $\rho_0:\Gamma\to\GL(W)$ one
of the retained $H_3$ representations, with $\dim W=3$. Write
$M_i=\rho_0(m_i)$ for the diagram meridians. Each $M_i$ is an
involution with eigenvalues $(-1,1,1)$ over
$\mathbb Q(\tau)$, where $\tau^2=\tau+1$.

A KCH deformation means a representation for which every
$\rho(m_i)-I$ has rank one and the nonidentity eigenvalue remains
different from one. For a formal such family, put
\[
 P(t)=\frac{\rho_t(m_0)-I}{\mu(t)-1},\qquad
 P_0=\frac{I-M_0}{2}.
\]
These are rank-one idempotents. The invertible matrix
$S=P(t)P_0+(I-P(t))(I-P_0)$ reduces to $I$ and satisfies
$P(t)S=SP_0$. Conjugation therefore normalizes the first meridian to
\begin{equation}\label{eq:normalize}
 M_0+s(t)(I-M_0),\qquad \mu(t)=-1+2s(t).
\end{equation}
No fixed-eigenvalue assumption is used in this normalization.

The retained four-seed Wirtinger propagation expresses all thirteen
arc matrices in terms of the four seed matrices. Use one variable
$s$ for the first seed and nine unrestricted infinitesimal entries
for each of the other three. There are $28$ variables. For a
crossing $M_b=M_oM_aM_o^{-1}$, its derivative is
\[
 \dot M_b=\dot M_oM_aM_o+M_o\dot M_aM_o-M_b\dot M_oM_o.
\]
For $M_b=M_o^{-1}M_aM_o$, it is
\[
 \dot M_b=-M_o\dot M_oM_b+M_o\dot M_aM_o+M_oM_a\dot M_o.
\]
Reverse propagation exchanges the two formulas. The actual crossing
signs are retained; involutivity of the base does not justify
discarding signs in the derivative.

Impose all thirteen crossing derivatives, giving $117$ scalar rows.
The rank-one tangent condition on another seed is equivalently
\[
 (I+M_i)\dot M_i(I+M_i)=0.
\]
Its nine entries give another $27$ rows. This condition is the
linearization of the rank-one minors, since $I+M_i$ projects, up
to the nonzero scalar two, onto the two-dimensional fixed space.
The resulting matrix is $144$ by $28$ over $\mathbb Q(\tau)$.

\begin{proposition}\label{prop:tangent}
For every retained base, that matrix has rank $23$ and nullity five.
Its kernel with $\dot s=0$ consists exactly of the four-dimensional
residual conjugation space. The quotient tangent space is
one-dimensional and $d\mu=2\,ds$ is nonzero on it.
\end{proposition}
The receipt supplies exact quadratic-rational kernel bases and
modular rank lower certificates in both split embeddings at
$11,19,29,31$. Exact independent kernel evaluation and basis
independence give the matching upper bound. The first-meridian
centralizer has dimension five; its scalar subspace acts trivially,
and its conjugation image has dimension four. Direct rational
calculations check this image lies in the kernel. A kernel vector
with $\dot s\ne0$ supplies the remaining direction.

The verifier rebuilds the crossing derivative blocks using
dual-number relator multiplication, with inverse derivative
$-M^{-1}\dot M M^{-1}$. It also tests every rank-one minor.
The PD parser and retained coefficient-group model are shared inputs;
the derivative formulas and symbolic nullspace solver are not shared.
The actual $T(3,5)$ family from the companion torus note gives a
positive tangent control, while adding its meridian-square relation
reduces the normalized nullity from five to four.

\section{Fixed meridian eigenvalue and the integrability gate}

\begin{corollary}\label{cor:fixed}
Every complex formal representation through one of these rank-three
points that keeps meridians involutive is formally conjugate to the
constant representation.
\end{corollary}
\begin{proof}
Formal involutions have the same eigenmultiplicities as their special
fibre, by their idempotents $(I\pm M)/2$. Thus the meridians still
satisfy the rank-one condition. Normalize the first as above, now
with $s=0$. Suppose all coefficients below order $j$ have been
removed by conjugation. At order $j$, the representation and
rank-one identities give an element of the fixed-$s$ tangent kernel
at $\rho_0$. By Proposition~\ref{prop:tangent}, it has the form
$[A,M_i]$ for a matrix $A$ commuting with $M_0$.
Conjugation by $I+t^jA$, with the sign chosen to remove this
coefficient, fixes $M_0$ and removes it on every seed, hence every
generator. The ordered product of these conjugations converges
$t$-adically. It conjugates the entire family to $\rho_0$.
\end{proof}

There is a precise remaining existence problem for varying $s$.
Choose four linear seed-coordinate conditions transverse to the
residual conjugations. The implicit function theorem applied to
the centralizer action puts all nearby deformations in this local
slice. With those four conditions added, the Jacobian has rank
$27$ in $28$ variables; its kernel is detected by $ds$.
Select $27$ independent equations. The analytic implicit function
theorem solves the other $27$ coordinates uniquely as convergent
functions of $s$. Substitute these functions in the remaining
equations. The result is a finite list of convergent functions
$f_j(s)$, all vanishing through first order.

Either every $f_j$ is identically zero, yielding an actual local
rank-three curve with parameter $s$, or some $f_j$ has a first
nonzero coefficient. In the latter case the local slice is an
isolated thick point, with analytic local algebra
$\C\{s\}/(s^e)$ for some $e\ge2$. This is a curve-or-obstruction
alternative, not a decision that the curve exists. A ramified
parameter change does not evade a nonzero $f_j$, because the
power-series ring of an arc is an integral domain.

\section{Restricted cocycles and a uniform stabilization theorem}

For a meridian-generated group with representation $\rho$ on $W$,
define
\[
 Z^1_{\rm adm}(\Gamma,W)=
 \{z\in Z^1(\Gamma,W):z(m_i)\in\im(\rho(m_i)-I)
 \text{ for every diagram meridian }m_i\}.
\]
Principal cocycles lie in this space. The admissibility condition
also holds on every conjugate meridian, by the cocycle identity.
Write $H^1_{\rm adm}=Z^1_{\rm adm}/B^1$ and use the analogous
definition for the contragredient module $W^*$.

\begin{theorem}[Formal splitting under restricted Ext vanishing]\label{thm:split}
Let $\Gamma$ be generated by meridians, and let
$\rho_0:\Gamma\to\GL(W)$ have rank-one meridian differences with
nonidentity eigenvalue $\mu_0\ne1$. Suppose
\[
 H^1_{\rm adm}(\Gamma,W)=0,
 \qquad H^1_{\rm adm}(\Gamma,W^*)=0.
\]
For any finite-dimensional trivial representation $T$, every
formal KCH family specializing to $\rho_0\oplus I_T$ is formally
conjugate, by a matrix reducing to $I$, to
\[
 \sigma_t\oplus I_T.
\]
The meridian eigenvalue of $\sigma_t$ is allowed to vary; no
constancy assumption is imposed. The statement holds over any field in which the
hypotheses hold.
\end{theorem}
\begin{proof}
Induct on the first order $j\ge1$ at which the off-diagonal blocks
or the lower-right deviation from $I_T$ have not yet been removed.
Write the blocks of a meridian, modulo $t^{j+1}$, as
\[
 \begin{pmatrix} A_i(t)&t^jB_i\\t^jC_i&I_T+t^jE_i\end{pmatrix},
 \qquad A_i(0)=\rho_0(m_i).
\]
A nonzero entry of $A_i(0)-I$ is a unit in the power-series ring.
The two-by-two rank-one minors, using that pivot entry, force
$E_i=0$, every column of $B_i$ to lie in
$\im(\rho_0(m_i)-I)$, and every row of $C_i$ to lie in the row
space of $\rho_0(m_i)-I$. Products of the two off-diagonal blocks
have order at least $2j>j$.

For a general group element, the first upper-right coefficient
satisfies
\[
 B(gh)=\rho_0(g)B(h)+B(g).
\]
Thus its columns are admissible cocycles. The lower-left
coefficient satisfies $C(gh)=C(g)\rho_0(h)+C(h)$.
The row $C(g)\rho_0(g)^{-1}$ is a cocycle for the row form of the
contragredient action, and its meridian rows are admissible.
The two vanishing hypotheses give constant matrices $U,V$ with
\[
 B(g)=(\rho_0(g)-I)U,\qquad
 C(g)=V(\rho_0(g)-I).
\]
Conjugation by
\[
 I+t^j\begin{pmatrix}0&-U\\V&0\end{pmatrix}
\]
removes both coefficients. It does not change the lower-right
coefficient at order $j$. That coefficient vanishes on the
meridians and is additive under products at this order, so it
vanishes throughout $\Gamma$. This advances the induction.
The product of the conjugating matrices converges $t$-adically.
In the limit the off-diagonal blocks vanish and the lower-right
block is $I_T$; the upper-left block is the required representation
$\sigma_t$. The rank-one condition is preserved by conjugation.
\end{proof}

\section{Application to all retained bases}

The restricted-cocycle matrix for $W$ has one scalar per meridian
and three rows per crossing, hence size $39$ by $13$. The same
construction is applied separately to $W^*$, whose meridian matrices
are $M_i^{-T}=M_i^T$. On all sixteen bases both matrices have rank
ten, so both restricted cocycle spaces have dimension three.
Direct principal-cocycle construction gives three independent
vectors in each. Consequently both restricted first cohomology
spaces vanish.

Exact quadratic-rational nullvectors and modular rank certificates
are replayed. A separate affine-cocycle evaluator checks the full
oriented knot relators and rebuilds every matrix column. It verifies
the three independent principal vectors directly. These are
characteristic-zero statements, covering both embeddings of
$\mathbb Q(\tau)$; no all-characteristic conclusion is asserted here.

\begin{corollary}
For every retained base and every $d\ge1$, a formal rank-one-meridian
family through $\rho_0\oplus I_d$ remains formally reducible.
It cannot supply an irreducible dimension-$3+d$ KCH lower bound.
\end{corollary}

In particular, an algebraic family with generic irreducible fibre
and this special fibre would give, after a local parameter change,
such a formal family, contradicting the splitting theorem. One can
see the contradiction by the full matrix-algebra criterion for
absolute irreducibility: a nonzero determinant of spanning word
matrices remains nonzero over the Laurent-series field, whereas a
split representation has no such determinant. This conclusion
concerns these special fibres only. A different component, another
semisimplification, a different meridian eigenvalue or an unrelated
rank-four representation remains possible.

\section{Finite jets and what remains unproved}

The frozen degree-three test uses both embeddings $\tau=4,8$
modulo $11$. Four transverse gauge coordinates and four local
rank-one minors per other seed are imposed together with every
crossing. An invertible selected Jacobian gives a unique candidate
jet with $s=t$. All $32$ candidates pass through degree three.
An independent replay uses entrywise polynomial adjugates, scalar
geometric-series inverses, all oriented relators and all rank-one
minors, rather than the producer's matrix-inverse recurrence or
four-minor subset.

These are residue-field jets, without a characteristic-zero lifting
or curve claim. A modular obstruction would additionally need a
verified transverse gauge and Jacobian over the local coefficient
ring. More modular success alone would not establish integrability.

The authoritative receipts are
\path{meridian_tangent_verification.json} and
\path{meridian_jet_verification.json}, with frozen inputs and exact
certificates retained. An actual curve requires proving the residual
identities. Sharp higher-rank representations require another special
fibre or mechanism. Full MRC and the sixteen target equalities stay open.
\end{document}
